\documentclass[10pt]{article}
\usepackage[margin=1in]{geometry}
\usepackage[T1]{fontenc}
\usepackage{lmodern,microtype,amsmath,amssymb,amsthm}
\theoremstyle{definition}
\newtheorem{assumption}{Assumption}
\begin{document}
\section*{Composite Minimization with Gap-Controlled Hessians}
We consider minimizing $F=f+\psi$ over a convex set $Q$ in a real Hilbert space $E$, where $f,\psi:E\to\mathbb{R}$ are convex on $Q$. At each feasible point, $f$ has a positive semidefinite Hessian satisfying $\|\nabla^2 f(x)\|_{\mathrm{op}}\leq H_0+H_1(F(x)-F(x^\star))$, where $H_0,H_1\geq0$ and $x^\star\in Q$ minimizes $F$. A given $x^0\in Q$ satisfies $\|x^0-x^\star\|\leq R$ and $F(x^0)-F(x^\star)\leq\Delta$, with $R,\Delta>0$. The oracle returns exact feasible gradients or exact proximal minimizers of $\psi$. For $\varepsilon>0$, the goal is $\bar{x}\in Q$ with $F(\bar{x})-F(x^\star)\leq\varepsilon$ almost surely.

The optimization problem is
\[
\min_{x\in Q} F(x),
\qquad F(x):=f(x)+\psi(x).
\]
The norm $\|\cdot\|$ is induced by the inner product on $E$, and $\|\cdot\|_{\mathrm{op}}$ denotes the induced operator norm. The algorithm receives $E$, $H_0$, $H_1$, $R$, $\Delta$, $\varepsilon$, and $x^0$; the functions and feasible set are accessed through the oracle.

Our assumptions:
\begin{assumption}[Convexity]
The set $Q\subseteq E$ is convex. For every $x,y\in Q$ and $t\in[0,1]$,
\[
\begin{aligned}
f((1-t)x+ty)&\leq(1-t)f(x)+tf(y),\\
\psi((1-t)x+ty)&\leq(1-t)\psi(x)+t\psi(y).
\end{aligned}
\]
\end{assumption}

\begin{assumption}[Minimizer and initialization]
There exists $x^\star\in Q$ such that $x^0\in Q$ and
\[
\begin{aligned}
F(x^\star)&\leq F(x) &&\text{for every }x\in Q,\\
\|x^0-x^\star\|&\leq R,\\
F(x^0)-F(x^\star)&\leq\Delta.
\end{aligned}
\]
\end{assumption}

\begin{assumption}[Feasible Hessian bound]
At every $x\in Q$, the function $f$ is Fr\'echet differentiable, and its gradient is Fr\'echet differentiable at $x$ as a map on $E$. The Hessian $\nabla^2 f(x):E\to E$ is a bounded self-adjoint linear operator satisfying
\[
\begin{aligned}
\langle\nabla^2 f(x)h,h\rangle&\geq0 &&\text{for every }h\in E,\\
\|\nabla^2 f(x)\|_{\mathrm{op}}
&\leq H_0+H_1\bigl(F(x)-F(x^\star)\bigr).
\end{aligned}
\]
These derivatives are ambient derivatives, including at boundary points of $Q$.
\end{assumption}

\begin{assumption}[Exact proximal access]
There is a map $\mathcal{P}:\mathbb{R}\times E\times E\to E$ such that, for every $a>0$ and $z,v\in E$,
\[
\mathcal{P}(a,z,v)\in
\operatorname*{arg\,min}_{u\in Q}
\left\{
\frac{\|u-z\|^2}{2a}
+\langle v,u-z\rangle+\psi(u)
\right\}.
\]
In particular, these minimizers exist and the returned point is feasible.
\end{assumption}

The oracle supports two query types. For a gradient query at $x\in Q$, the oracle returns $\nabla f(x)$ at $x$. For a proximal query $(a,z,v)\in\mathbb{R}\times E\times E$, it returns $\mathcal{P}(a,z,v)$, with the minimization property above applying when $a>0$. Both query types count toward one common oracle-call budget.

After the allotted calls, the output must satisfy
\[
\bar{x}\in Q,
\qquad F(\bar{x})-F(x^\star)\leq\varepsilon
\quad P\text{-a.s.},
\]
where $P$ is the joint law of the algorithm's and oracle's randomness, which are independent.
\end{document}
